\documentclass{article}

\usepackage[T1]{fontenc}

\usepackage{booktabs,graphicx,siunitx}

\begin{document}

\section{Poisson / negative binomial}

\begin{tabular}{lll}
\toprule
 & Log-count (B) & IRR \\
\midrule
(Intercept) & $-$1.09 & 0.34 \\
 & (0.06) & [0.30, 0.38] \\
teaching\_method: Flipped & $-$0.56 & 0.57 \\
 & (0.08) & [0.49, 0.67] \\
teaching\_method: Lecture & 0.28 & 1.33 \\
 & (0.06) & [1.17, 1.51] \\
gender: Male & $-$0.04 & 0.96 \\
 & (0.06) & [0.86, 1.07] \\
\midrule
Num.Obs. & 240 &  \\
Dispersion & 3.14 &  \\
Residual deviance & 633.11 &  \\
df & 236 &  \\
Log.Lik. & $-$683.51 &  \\
AIC & 1375.03 &  \\
BIC & 1388.95 &  \\
\bottomrule
\end{tabular}

\includegraphics{figures/01_poisson-negative-binomial/figure_residuals.png}

Models event counts. Each predictor's incidence-rate ratio (IRR) gives the multiplicative change in the expected count per unit ($>$1 more, $<$1 fewer); p tests significance. Check dispersion to pick Poisson vs. negative binomial. If cases have unequal exposure (different observation time/area/population), add an offset of log(exposure) so the model predicts rates, not raw counts - omitting it biases the IRRs. Your significance threshold ($\alpha$) is 0.05.

Predictor teaching\_method: Flipped was associated with the count, IRR=0.57, 95\% CI [0.49, 0.67], p $<$ .001.

Statistical basis: Poisson regression (incidence-rate ratios) (R Core Team (2026). R: A Language and Environment for Statistical Computing. R Foundation for Statistical Computing, Vienna.); Negative binomial regression (glm.nb) for overdispersed counts (Venables WN, Ripley BD (2002). Modern Applied Statistics with S, Fourth edition. Springer, New York.); Overdispersion check (Lüdecke D, Ben-Shachar MS, Patil I, Waggoner P, Makowski D (2021). "performance: An R Package for Assessment, Comparison and Testing of Statistical Models." Journal of Open Source Software, 6(60), 3139.).

\end{document}